\documentclass[oneside,final,14pt,a4paper]{extreport}

\usepackage{tempora}

\usepackage{vmargin}
\setpapersize{A4}
\setmarginsrb{2.5cm}{2.0cm}{2.0cm}{2.0cm}{0pt}{10mm}{0pt}{13mm}
\usepackage{setspace}
\setstretch{1.5}
\usepackage{indentfirst}
\parindent=1.25cm

\DeclareFontShape{OT1}{cmtt}{bx}{n}{<5><6><7><8><9><10><10.95><12><14.4><17.28><20.74><24.88>cmttb10}{}
\renewcommand{\ttdefault}{pcr}
\renewcommand{\thesection}{\arabic{section}}

\usepackage{atbegshi,picture}
\usepackage[T1,T2A]{fontenc}
\usepackage[utf8]{inputenc}

\usepackage[russian]{babel}
\usepackage[backend=biber,style=ieee,sorting=none,autocite=inline]{biblatex}
\bibliography{ref.bib}
\DefineBibliographyStrings{english}{%
  bibliography = {References},}

\usepackage{pdfpages}
\usepackage{url}
\usepackage{hyperref}
\hypersetup{hidelinks}
\usepackage{enumitem}
\usepackage{fancyhdr}

\pagestyle{fancyplain}

\renewcommand{\chaptermark}[1]%
	{\markboth{\chaptername~\thechapter~--~#1}{}}

\renewcommand{\sectionmark}[1]%
	{\markright{\thesection\ #1}}

\rhead[\fancyplain{}{\bf\leftmark}]%
      {\fancyplain{}{\bf\thepage}}
\lhead[\fancyplain{}{\bf\thepage}]%
      {\fancyplain{}{\bf\rightmark}}
\cfoot{}

\begin{document}

\includepdf[pages=1, offset=70 -80]{title.pdf}
\tableofcontents
\newpage

\setcounter{page}{3}
\setcounter{section}{0}

\section{Введение}

Выпускная квалификационная работа посвящена проектированию, реализации и оценке среды выполнения контейнеров OCI, ориентированной на безопасность. Актуальность темы определяется тем, что контейнеры разделяют ядро хоста, тогда как эталонный \texttt{runc} ориентирован на производительность, а не на изоляцию~\cite{wang2022,volpert2024}. Seccomp, AppArmor и capabilities требуют ручной настройки; песочницы gVisor и Kata дают сильную изоляцию ценой 20--100\% накладных расходов~\cite{viktorsson2020}; кластерные рекордеры профилей предполагают инфраструктуру оркестрации. Операторы plain \texttt{runc} остаются без интегрированного инструмента автоматизации.

Цель --- спроектировать и экспериментально оценить форк \texttt{runc}, генерирующий профили seccomp, AppArmor и capabilities из наблюдаемого поведения нагрузки одним флагом CLI. Объектом исследования является изоляция контейнерных процессов на Linux в рамках OCI. Предметом --- интегрированная runtime-локальная генерация профилей confinement. Научная новизна заключается в объединении трёх механизмов ядра в двухфазный конвейер «сканирование --- принудительное исполнение» с cgroup-wide трассировкой capabilities и самовыполняемыми OCI-хуками внутри бинарного файла runtime.

\section{Краткое содержание}

В основной части выполнен обзор литературы по контейнерной безопасности, изолированным runtime и автоматической генерации seccomp~\cite{lopes2020}. Сравнения runtime, как правило, противопоставляют \texttt{runc} с политиками по умолчанию песочницам, не измеряя потенциал hardened shared-kernel runtime; автоматизация чаще охватывает только seccomp, а AppArmor и capabilities остаются фрагментированными~\cite{mattetti2020}.

Проектная часть формулирует требования к сканеру: однофлаговый UX, три профиля за один проход, совместимость с upstream \texttt{runc}, минимальные зависимости и финализация, которая никогда не расширяет привилегии. Модель угроз охватывает эскалацию через неиспользуемые syscalls и capabilities, выход за пределы файловой системы и post-exploitation после компрометации приложения; уязвимости прикладного уровня и supply-chain атаки вынесены за рамки.

Независимо от режима сканирования форк меняет дефолт безопасности: \texttt{runc spec} теперь генерирует пустой набор capabilities, а исторические наборы upstream (3 capabilities) и Docker (14 capabilities) доступны через два взаимоисключающих флага. Архитектура реализует режим \texttt{--security-scan}: ослабление confinement в памяти, установка хуков, запуск контейнера, остановка трассировщиков, сужение \texttt{config.json}. Seccomp записывается через runtime-трассировку syscalls; capabilities --- через eBPF-трассировку по cgroup~v2; AppArmor обучается в complain mode. Хуки реализованы подкомандами \texttt{runc}; OCI bundle не получает исполняемых артефактов.

Реализация --- форк \texttt{opencontainers/runc}. Лабораторный стенд \texttt{runc-hardened-test} на Ubuntu~24.04 сравнивает stock \texttt{runc} и hardened runtime на уязвимом Flask-приложении: baseline, scan с легитимным трафиком, enforcement и повтор post-exploitation payload. Один проход \texttt{--security-scan} формирует \texttt{seccomp.json}, профиль AppArmor и журнал capabilities. После scan-then-enforce reconnaissance, чтение \texttt{/etc/shadow}, reverse shell и mount probe блокируются enforced seccomp; ping и upload календаря сохраняются после документированных корректировок. Ограничения: покрытие равно пройденным путям; file capabilities не попадают в trace; сравнение с gVisor и Kata носит аналитический характер по данным~\cite{wang2022,volpert2024}.

\section{Заключение}

В результате работы разработан и экспериментально оценён security-first форк OCI runtime \texttt{runc}. Форк вносит три изменения в слой runtime: пустой набор capabilities по умолчанию с двумя opt-in флагами, флаг \texttt{--security-scan} для генерации профилей seccomp, AppArmor и capabilities, и сужение \texttt{config.json} до наблюдаемого набора привилегий без его расширения. Основной вклад --- интегрированная генерация профилей seccomp, AppArmor и capabilities в воспроизводимом runtime-локальном конвейере с одним флагом \texttt{--security-scan}. Полученные результаты подтверждают существенное улучшение post-exploitation containment относительно stock \texttt{runc} на тестируемой нагрузке при сохранении легитимного поведения приложения. Выводы ограничены одним web-workload и качественными метриками безопасности; дальнейшая работа включает слияние file capabilities, расширение лаборатории и количественную оценку производительности относительно изолированных runtime~\cite{reeves2021runtime}.

\section{Список использованной литературы}

{\small
\setlength{\parskip}{0.15em}
\begin{enumerate}[label=\arabic*., leftmargin=1.5em, itemsep=0.15em, parsep=0pt]
\item Wang X., Johannesson D. Performance and isolation analysis of RunC, gVisor and Kata Containers runtimes // Cluster Computing. 2022. Vol.~25. P.~2363--2380.
\item Volpert S., Winkelhofer S., Wesner S., Domaschka J. An Empirical Analysis of Common OCI Runtimes' Performance Isolation Capabilities // Proceedings of the 15th ACM/SPEC ICPE. 2024. P.~60--70.
\item Viktorsson W., Klein C., Tordsson J. Security-Performance Trade-offs of Kubernetes Container Runtimes // 2020 IEEE MASCOTS. 2020. P.~1--4.
\item Lopes N., Martins R., Correia M. E., Serrano S., Nunes F. Container Hardening Through Automated Seccomp Profiling // Proceedings of the 6th International Workshop on Container Technologies and Container Clouds. 2020. P.~31--36.
\item Mattetti M., Allouche Y. Automatic security hardening and protection of linux containers // Computers \& Security. 2020. Vol.~96. Article 101921.
\item Reeves M., Tian D. J., Bianchi A., Celik Z. B. Towards Improving Container Security by Preventing Runtime Escapes // 2021 IEEE SecDev. 2021. P.~38--46.
\end{enumerate}
}

\end{document}
